\section{Exact operator bridge to the released Navier--Stokes profiles}
\label{nsbridge:nsb:operator}

This bounded comparison uses the locally saved released Navier--Stokes source
(the source pages are 7--15 and 24--27 for the profile, stress and pulse
construction). Its stated physical viscosity is denoted by $\nu_{\rm NS}>0$.
The Boussinesq lane variables and physical viscosity remain $d$, $\mu$ and
$\nu=\sqrt{A_0}$; none is identified with $\nu_{\rm NS}$.

Use right-handed cylindrical coordinates $(r,\vartheta,z)$ and define
\[
 y_1=z,\qquad y_2=\frac{r^2}{2}>0,\qquad
 v=(u^z,ru^r)=J\nabla_y\Psi,
 \quad J=\begin{pmatrix}0&-1\\1&0\end{pmatrix}.
\]
The meridional divergence identity is
$\partial_{y_1}v_1+\partial_{y_2}v_2=0$. Define the circulation and
meridional vorticity variables
\[
 \Gamma=ru^\vartheta,
 \qquad \xi=\frac{\omega^\vartheta}{r},
 \qquad h=\frac{\Gamma}{2y_2}=\frac{u^\vartheta}{r},
 \qquad Q=\partial_{y_1}^2+2y_2\partial_{y_2}^2,
 \qquad M=Q+4\partial_{y_2}.
\]
The azimuthal momentum and meridional vorticity rows of the cylindrical
equations (with arbitrary smooth force components) become, on $y_2>0$,
\begin{align}
 D_v\Gamma&=\nu_{\rm NS}Q\Gamma+r f^\vartheta,
 \label{nsbridge:nsb:Gamma}\\
 D_v\xi&=\nu_{\rm NS}M\xi+\partial_{y_1}(h^2)
       +\operatorname{curl}_y(f^z,f^r/r),
 \label{nsbridge:nsb:xi}
\end{align}
where $D_v=\partial_t+v\cdot\nabla_y$ and
$\operatorname{curl}_y(a,b)=\partial_{y_1}b-\partial_{y_2}a$.

Indeed, $\partial_r=r\partial_{y_2}$ and
$\partial_r^2+r^{-1}\partial_r=2y_2\partial_{y_2}^2+2\partial_{y_2}$.
For $\Gamma$, the azimuthal viscous operator is
$\partial_r^2-r^{-1}\partial_r+\partial_z^2=Q$.
For $r\xi=\omega^\vartheta$, the cylindrical operator
$(\partial_r^2+r^{-1}\partial_r-r^{-2}+\partial_z^2)(r\xi)$,
divided by $r$, is $M\xi$; explicitly the radial terms are
$2y_2\xi_{y_2y_2}+4\xi_{y_2}$. The centrifugal term is
$r^{-2}\partial_z((u^\vartheta)^2)=\partial_{y_1}(\Gamma^2/(4y_2^2))
=\partial_{y_1}(h^2)$. These calculations prove every coefficient in
(\ref{nsbridge:nsb:Gamma})--(\ref{nsbridge:nsb:xi}), including the $r$ force conversion.

The velocity--vorticity relation is also retained. Since
$u^z=-\Psi_{y_2}$ and $u^r=\Psi_{y_1}/r$,
\[
 \xi=\frac{\partial_z u^r-\partial_r u^z}{r}
     =\frac{\Psi_{y_1y_1}}{2y_2}+\Psi_{y_2y_2}
     =\frac{Q\Psi}{2y_2}.
\]
Thus the two transport rows do not by themselves replace the elliptic
velocity reconstruction. Conversely, on a simply connected meridional
domain, a smooth $\Psi,\Gamma$ satisfying this relation and both transport
rows recovers $u$. Define its momentum residual without pressure as
$a=\partial_tu+(u\cdot\nabla)u-\nu_{\rm NS}\Delta u-f$.
The circulation row is $a^\vartheta=0$, while the vorticity row is
$\partial_z a^r-\partial_r a^z=0$. The one-form
$a^r\,\mathrm dr+a^z\,\mathrm dz$ is therefore closed, and
$p=-\int(a^r\,\mathrm dr+a^z\,\mathrm dz)$ along a path from a fixed
base point is path independent. Its derivatives are $p_r=-a^r$,
$p_z=-a^z$. This proves the local momentum reconstruction, with pressure
unique up to a function of time; boundary data are not removed.

The change from circulation to swirl rate is invertible on the exact retained
domain: $\Gamma=2y_2h$. Product differentiation gives
\[
 Q(2y_2h)=2y_2Mh,\qquad
 D_v(2y_2h)=2y_2D_vh+2v_2h.
\]
Therefore (\ref{nsbridge:nsb:Gamma}) is equivalent to
\begin{equation}
 D_vh=\nu_{\rm NS}Mh-\frac{v_2}{y_2}h+\frac{f^\vartheta}{r}.
 \label{nsbridge:nsb:h}
\end{equation}
The source profile uses precisely this relation: in its similarity variables,
$E$ is the azimuthal velocity profile and
$H_{\rm src}=\sqrt{2X}E$ is the scaled circulation profile (source equations
(4.3), (4.8)). Since $r=\sqrt{2qX}$ and
$u^\vartheta=q^{-A}E$, its physical circulation is
$\Gamma=q^{1/2-A}H_{\rm src}=q^{-h_{\rm src}}H_{\rm src}$,
where the source exponents satisfy $A=1/2+h_{\rm src}$. Here $q>0$ is the
geometric similarity coordinate and $h_{\rm src}>0$ is an exponent; neither
is a signed amplitude scale. This is the exact scale dictionary; it does not
identify $H_{\rm src}$ with the Boussinesq scalar amplitude.

The sign of the blowout amplitude must be kept separate from that positive
geometric coordinate. The axisymmetric equations have the exact involution
\[
 (u^r,u^\vartheta,u^z,p;f^r,f^\vartheta,f^z)
 \longmapsto
 (u^r,-u^\vartheta,u^z,p;f^r,-f^\vartheta,f^z).
 \label{nsbridge:nsb:sign}
\]
It sends $(\Gamma,h)$ to $(-\Gamma,-h)$, leaves $\xi$ and the meridional
flow unchanged, and preserves every displayed equation: the only quadratic
swirl term is $\partial_{y_1}(h^2)$, while the azimuthal equation is linear in
$(\Gamma,f^\vartheta)$. Consequently a source-selected positive profile
$E>0$ (source pp. 7--9) does not justify assuming that every blowout branch is
positive; its reflected branch has $E<0$, $\Gamma<0$ and the same blowup
magnitude. The exponent $q^{-h_{\rm src}}$ controls the magnitude in either
branch. The source's positive covariance coefficients and viscosity are
separate construction choices, not a sign claim about the velocity amplitude.

For comparison with the lane's controlled scalar equation, set
$B=h^2$. The globally smooth square equation, including its geometric and
viscous terms, is
\begin{equation}
 D_vB=\nu_{\rm NS}MB-2\nu_{\rm NS}
 |\nabla h|_Q^2-\frac{2v_2}{y_2}B+\frac{2h}{r}f^\vartheta,
 \qquad |\nabla h|_Q^2=h_{y_1}^2+2y_2h_{y_2}^2.
 \label{nsbridge:nsb:B}
\end{equation}
This follows by multiplying (\ref{nsbridge:nsb:h}) by $2h$ and using
$M(h^2)=2hMh+2|\nabla h|_Q^2$. At zeros of $h$ equation
(\ref{nsbridge:nsb:B}) remains smooth; dividing by $B$ is unnecessary and would lose
that domain. On a component where $B>0$, the equivalent expression is
$D_vB=\nu_{\rm NS}MB-\nu_{\rm NS}|\nabla B|_Q^2/(2B)
-2v_2B/y_2+2hf^\vartheta/r$.

The exact relation to the Boussinesq lane is consequently an operator
correspondence, not a scalar conjugacy. The lane's retained-viscosity scalar
row has $D_v\Theta=d\Delta\Theta+f_\Theta$ in Cartesian material
coordinates. A direct identification $\Theta=B$ would require, on the same
physical domain, both $M=\Delta$ after the nonlinear coordinate pullback and
$-2v_2B/y_2-2\nu_{\rm NS}|\nabla h|_Q^2+2hf^\vartheta/r$ to be absorbed into
one prescribed force. The first condition fails as an operator identity:
$M=\partial_{y_1}^2+2y_2\partial_{y_2}^2+4\partial_{y_2}$ has variable metric
and drift, whereas the lane's $\Delta$ is the Euclidean Cartesian Laplacian.
The geometric, gradient and force terms must therefore be retained or
calculated in the particular source state. Nonconstant profiles alone
would not prove that such terms never cancel.
The exact bridge is (\ref{nsbridge:nsb:Gamma})--(\ref{nsbridge:nsb:B}) with the
full residual terms retained. It neither claims that the lane's controls are
the source pulses nor changes the source's pressure reconstruction.

For the meridional velocity, the source's pulse construction supplies the
following further exact local dictionary (source pp. 11--15 and 74--77). If
$w$ is a divergence-free pulse and $\pi$ its pressure increment, then
\[
 R(u_B+w,p_B+\pi)=R(u_B,p_B)+L_{u_B}(w,\pi)
 +\nabla\!\cdot(w\otimes w),
\]
\[
 L_{u_B}(w,\pi)=\partial_tw+(u_B\cdot\nabla)w
 +(w\cdot\nabla)u_B-\nu_{\rm NS}\Delta w+\nabla\pi,
 \qquad \nabla\cdot w=0.
\]
For the moving local phase with wavevector $n$ and transverse amplitude
$t_m$, the source equation is
\[
 t_m'+\mathcal Kt_m+m^2d_N t_m+i k m n\,\pi_m=-f_m,
 \qquad n\cdot t_m=0,
 \qquad d_N=\varepsilon k^2|n|^2,
\]
with the source's moving shear matrix
$\mathcal K=\begin{pmatrix}0&-2F&0\\2F+RF_R&0&0\\G_R&0&0\end{pmatrix}$.
The transverse projection uses
\[
 A_\Phi=-\mathcal K+
 \frac{n(n^T\mathcal K-n'^T)}{|n|^2}.
\]
These are exact source pulse equations, including viscosity and pressure
projection. Their quadratic covariance is chosen to cancel the annular
stress; the source explicitly retains curl, cutoff, and correction interactions.
The Boussinesq lane's two-dimensional pair and its full profile residuals have
an exact algebraic form, but no identification with this three-dimensional
pulse equation follows without an additional embedding that supplies the
source's cylindrical geometry, moving plane, covariance cone and correction
cycle. That embedding is not asserted here.

The current bridge therefore proves a precise morphism of variables, operators,
forces and pulse residuals and records the exact obstruction to a direct scalar
identification. It establishes the next programme datum while preserving the
unproved status of any singular limit or third return.


\sloppy

\section{Scope and source equations}
This note records one local, exact comparison.  It does not identify the
coupled Boussinesq stages with the released Navier--Stokes construction and
makes no claim about a global Navier--Stokes solution.  The source is the
released PDF \texttt{openai\_navier\_stokes.pdf}, SHA-256
\texttt{8c8a94ad9ac824c8b605b9827cadf7beaca48bd10b380de3cfc872a2c37afa81}.
The source statements used here are Section 7, equations (7.2)--(7.8),
(7.12)--(7.22), and Section 9, equations (9.1)--(9.2), on PDF pages 74--75,
77--81, and 101--103. In the source notation, on one lifted pulse rectangle,
\begin{equation}
 \Phi=p\theta+p_zZ/\varepsilon+x_0R-v(pF+p_zG),\qquad
 n=\nabla_*\Phi,
\end{equation}
with
\begin{equation}
 n=(x_0-v(pF_R+p_zG_R),\ p/R,\ p_z-\varepsilon v(pF_Z+p_zG_Z)).
 \label{nsbridge:eq:nsource}
\end{equation}
For a nonzero integer harmonic $m$, the exact principal pulse equations are
\begin{equation}
 n\cdot t_m=0,\qquad
 t_m'+\mathsf Kt_m+m^2d_Nt_m+ikm n\pi_m=-f_m,
 \qquad d_N=\varepsilon k^2|n|^2,
 \label{nsbridge:eq:sourcepulse}
\end{equation}
where a prime is the pulse-coordinate derivative and
\begin{equation}
 \mathsf K=\begin{pmatrix}0&-2F&0\\2F+RF_R&0&0\\G_R&0&0\end{pmatrix},\qquad
 \mathsf A_\Phi=-\mathsf K+\frac{n(n^T\mathsf K-(n')^T)}{|n|^2}.
 \label{nsbridge:eq:sourceoperator}
\end{equation}
The source pressure coefficient is exactly
\begin{equation}
 \pi_m=\frac{i}{km|n|^2}\bigl(n\cdot\mathsf Kt_m-(n')\cdot t_m+n\cdot f_m\bigr).
 \label{nsbridge:eq:sourcepressure}
\end{equation}
For $m=1$, a real homogeneous pulse $t_1$ has the source energy identity
(7.22):
\begin{equation}
 \frac12(|t_1|^2)'=-g\cdot\bigl(t_{1,r}(t_{1,\theta},t_{1,z})\bigr)
 -\varepsilon k^2|n|^2|t_1|^2,
 \label{nsbridge:eq:sourceenergy}
\end{equation}
where $g=(RF_R,G_R)$ in the tangential coordinates.  Thus the source pulse
retains both shear transfer and viscous dissipation.

\section{The coupled physical pair}
For one activated layer of the coupled calculation, put
\begin{equation}
 a_c=-\frac{J\zeta\cdot G_{<q}}{Kr^2},\qquad b_c=K\zeta_1,\qquad
 \lambda_c^\Theta=\kappa K^2r^2,\qquad \lambda_c^\Omega=\varepsilon_c K^2r^2.
\end{equation}
Its exact retained-viscosity pair, including optional scalar and vorticity
controls, is
\begin{equation}
 \dot y=C_cy+c_c,\qquad y=\binom{\Theta}{\Omega},\qquad
 C_c=\begin{pmatrix}-\lambda_c^\Theta&a_c\\b_c&-\lambda_c^\Omega\end{pmatrix},\quad
 c_c=\binom{c_\theta}{c_\omega}.
 \label{nsbridge:eq:coupledpair}
\end{equation}
Here $G_{<q}$, $\zeta$, and $r$ evolve with the inherited background and phase
ODEs; in particular no fixed-background substitution is made.  Let $v$ be a
source pulse coordinate on an interval and let $q_v=\dd t/\dd v$ describe
its physical-time parametrization. The forward source parametrization has
$q_v>0$; the algebra below also retains a negative orientation without
changing the amplitude sign. Define $y'(v)=q_v(C_cy+c_c)$.
Here $K,r$ are the coupled frequency and phase length, not the source
carrier $k$ and cylindrical radius $R$. The original $A_0,\lambda_0,\mu$
and time scale $\nu=\sqrt{A_0}$ are unchanged. The source chart parameter
$\varepsilon$ is not identified with either physical diffusivity.

\section{Exact local residual map}
Let $B(v)\in\R^{3\times2}$ be a $C^1$ tangent frame with full column rank and
$n(v)^TB(v)=0$, so its range is exactly $n(v)^\perp$.  Let $B_\ell(v)$ be any
left inverse, $B_\ell B=I_2$; then $B B_\ell$ acts as the identity on
$n^\perp$.  Let $L(v)\in\R^{2\times2}$ be any $C^1$ coefficient map and embed
the coupled pair into a transverse source amplitude by
\begin{equation}
 t(v)=B(v)L(v)y(v).
 \label{nsbridge:eq:embedding}
\end{equation}
For an embedding with a non-tangent frame the transversality condition would be
\begin{equation}
 n(v)^TB(v)L(v)y(v)=0.
 \label{nsbridge:eq:transversecondition}
\end{equation}
For the tangent frame used here, the condition holds for every $L,y$.  Set
\begin{equation}
 H(v)=B_\ell(v)\bigl(\mathsf A_\Phi(v)B(v)-B'(v)\bigr).
 \label{nsbridge:eq:Hdef}
\end{equation}
Then the exact projected source residual of the embedded coupled pair is
\begin{align}
 \mathcal E_{L,c}:=B_\ell\proj_{n^\perp}
 \left[t'-\mathsf A_\Phi t+m^2d_Nt+f_m\right]
 &=\left[L'+q_vLC_c-HL+m^2d_NL\right]y+q_vLc_c \\
 &\quad+B_\ell\proj_{n^\perp}f_m.
 \label{nsbridge:eq:residualmap}
\end{align}
In particular, with $L=I$ and no added source, the defect is the explicit
matrix expression
\begin{equation}
 \mathcal E_{I,0}=\bigl(q_vC_c-H+m^2d_NI_2\bigr)y.
 \label{nsbridge:eq:identitydefect}
\end{equation}
This gives a local residual relation with every actual coefficient retained.

\begin{proof}
Differentiate \eqref{nsbridge:eq:embedding}:
$t'=B'L y+B L'y+B L y'$.  Since $y'=q_v(C_cy+c_c)$,
differentiate $n^TB=0$ to obtain $n^TB'=-(n')^TB$. Also
$n^T\mathsf A_\Phi B=-(n')^TB$, directly from
\eqref{nsbridge:eq:sourceoperator}. Thus $\mathsf A_\Phi B-B'$ has transverse
range and equals $BH$. Applying $B_\ell$ to the transverse terms and
$B_\ell\operatorname{proj}_{n^\perp}$ to the force proves
\eqref{nsbridge:eq:residualmap}. More explicitly,
\begin{equation}
 n^T(t'-\mathsf A_\Phi t+m^2d_Nt)=0,\qquad
 n^T(t'-\mathsf A_\Phi t+m^2d_Nt+f_m)=n^Tf_m.
\end{equation}
The unprojected force need not be transverse. The pressure below removes
precisely this remaining normal component.
\end{proof}

\section{Pressure reconstruction and full three-dimensional residual}
For any embedded transverse $t$ and arbitrary $f_m$, define
\begin{equation}
 \pi_{L,c}=\frac{i}{km|n|^2}
 \left(n\cdot\mathsf Kt-(n')\cdot t+n\cdot f_m\right).
 \label{nsbridge:eq:pressuremap}
\end{equation}
Then the full source residual satisfies the exact decomposition
\begin{align}
 \mathcal R_{L,c}
 &: =t'+\mathsf Kt+m^2d_Nt+ikm n\pi_{L,c}+f_m \\
 &=\proj_{n^\perp}\left(t'+\mathsf Kt+m^2d_Nt+f_m\right)
 =B\,\mathcal E_{L,c},
 \label{nsbridge:eq:fullresidual}
\end{align}
where the last equality uses the definition of $H$ and the transverse condition.
Thus the pressure removes precisely the normal component; it does not remove
$\mathcal E_{L,c}$, the moving-frame, shear, diffusion, or coupled-background
mismatch.  If one elects to force the source pulse to obey the released equation,
then the unique additional transverse forcing is
\begin{equation}
 f_{m,\mathrm{corr}}=-B\,\mathcal E_{L,c},
\end{equation}
added to the old $f_m$. The old $\mathcal E_{L,c}$ is evaluated before
this addition, so the new projected residual is zero.
Since $n^Tf_{m,\mathrm{corr}}=0$, pressure is unchanged.
This is an exact local
forced relation, not an identification of the physical constructions.

\section{Diffusion and energy comparison}
The source diffusion multiplier is $m^2d_N=m^2\varepsilon k^2|n|^2$ in
\eqref{nsbridge:eq:sourcepulse}.  The coupled pair has the two retained physical losses
$\lambda_c^\Theta$ and $\lambda_c^\Omega$, which are generally unequal to
$m^2d_N$.  Their signed discrepancy in \eqref{nsbridge:eq:identitydefect} is explicitly
\begin{equation}
 \bigl(m^2d_N-q_v\lambda_c^\Theta\bigr)\Theta,
 \qquad
 \bigl(m^2d_N-q_v\lambda_c^\Omega\bigr)\Omega
\end{equation}
when $H=0$ and $C_c$ is expanded entrywise; with general $H$ the exact
diagonal entries are those of $q_vC_c-H+m^2d_NI_2$.  No damping is dropped.
If one compares nonnegative loss magnitudes instead, the corresponding cost
gaps are $|m^2d_N-q_v\lambda_c^\Theta|$ and
$|m^2d_N-q_v\lambda_c^\Omega|$.
The required finite-interval coordinate map can be constructed explicitly.
Let $I=[v_0,v_1]$, and let $X_c,X_s$ be the identity-initial-value
fundamental matrices for $D_c=q_vC_c$ and $D_s=H-m^2d_NI_2$,
respectively. For either continuous matrix $D$, define
\[
 X_D(v)=I_2+\sum_{j\ge1}
 \int_{v_0\le w_j\le\cdots\le w_1\le v}
 D(w_1)\cdots D(w_j)\,\dd w_j\cdots\dd w_1 .
\]
With $M_D=\sup_I\|D\|$ and $T_I=v_1-v_0$, the $j$th integral has norm
at most $(M_DT_I)^j/j!$. Uniform convergence proves the integral equation
$X_D=I_2+\int_{v_0}^vDX_D$ and hence $X_D'=DX_D$.
The same iteration for $Z'=-ZD$, $Z(v_0)=I_2$ gives $ZX_D=I_2$
by differentiation, proving invertibility on $I$. For a specified
invertible $L_0$, set
\[
 L=X_sL_0X_c^{-1},\qquad L^{-1}=X_cL_0^{-1}X_s^{-1}.
\]
This has norm at most $\|L_0\|e^{(M_c+M_s)T_I}$; its inverse has the
same bound with $L_0^{-1}$. Differentiation proves
\begin{equation}
 L'=HL-q_vLC_c-m^2d_NL,
 \label{nsbridge:eq:intertwiner}
\end{equation}
so the homogeneous embedded pair has zero projected residual. This proves
an actual finite-interval conjugacy along the retained coefficient path.
It does not make the source phase,
pressure, cutoff, or nonlinear covariance equal to the coupled Boussinesq fields.
For general controls $c_c$, the forced term is $q_vLc_c$ in
\eqref{nsbridge:eq:residualmap}.

For an arbitrary embedded field (possibly complex), the exact energy balance
retains its full defect:
\begin{equation}
 \begin{split}
 \frac12(|t|^2)'={}&-\operatorname{Re}(\overline t^{\,T}\mathsf Kt)
 -m^2d_N|t|^2-\operatorname{Re}(\overline t^{\,T}f_m)\\
 &+\operatorname{Re}(\overline t^{\,T}B\mathcal E_{L,c}).
 \end{split}
\end{equation}
This follows by taking the real Hermitian product of
\eqref{nsbridge:eq:fullresidual} with $t$. Its pressure term vanishes because $n$ is
real and $n^Tt=0$. The $2F$ part of $\mathsf K$ is skew-symmetric, hence
its real work is zero; the shear work equals
$\operatorname{Re}(\overline{t_r}(g_\theta t_\theta+g_z t_z))$.
Only for a zero residual, zero force, $m=1$, and real $t$ does this reduce
to \eqref{nsbridge:eq:sourceenergy}. The real coupled pair has
\begin{equation}
 \frac12\frac{\dd}{\dd t}|y|^2
 =-\lambda_c^\Theta\Theta^2-\lambda_c^\Omega\Omega^2
 +(a_c+b_c)\Theta\Omega+c_\theta\Theta+c_\omega\Omega,
 \label{nsbridge:eq:coupledenergy}
\end{equation}
which is a two-component physical energy identity.  Equations
\eqref{nsbridge:eq:residualmap} and \eqref{nsbridge:eq:coupledenergy} exhibit exactly where the
source's three-dimensional shear flux and the coupled pair's two-dimensional
exchange differ.

\section{Proven scope}
The source uses a cylindrical, localized, divergence-free curl construction
(Section 7.4 and Section 9, PDF pages 85--87 and 102--103); the displayed
principal pulse equation is only one component of that construction.  The
comparison above proves an exact local map and its full pressure-projected
residual for every finite interval on which the denominators and frame left
inverse exist.  It preserves source viscosity, moving phase normal, pressure,
rotation/shear matrix, and all coupled physical diffusion coefficients.  It does
not prove that the coupled third-growth or steering stages produce the source's
pulse phase $\Phi$, covariance, cutoff tails, or smooth terminal force.  Those
remain separate bridge problems.

The sign conventions in this comparison are also explicit.  The quantities
$q_v>0$, $\varepsilon>0$, and $d_N=\varepsilon k^2|n|^2>0$ record the chosen
orientation of physical time and the positive viscous coefficient; they are
not assumptions on the sign of a pulse or blowout amplitude.  The equations
\eqref{nsbridge:eq:sourcepulse}--\eqref{nsbridge:eq:sourcepressure} are linear in $(t_m,f_m)$,
so the simultaneous replacement $(t_m,f_m,\pi_m)\mapsto
(-t_m,-f_m,-\pi_m)$ is an exact reflected branch.  Likewise, the local map
\eqref{nsbridge:eq:embedding} is sign-equivariant under $y\mapsto-y$ and
$c_c\mapsto-c_c$, with the coefficients fixed at the actual path.
This is a statement about the retained linear equations, not a sign symmetry
of the entire coupled nonlinear system.
The source covariance is quadratic, so reversing a wave leaves that
covariance unchanged. Reversing the full Navier--Stokes velocity need not
reverse its force, since $(u\cdot\nabla)u$ is even in $u$.
The full-flow map that reverses axisymmetric swirl is a spatial reflection;
its complete formula is proved in the accompanying covariance bridge.
No sign of a blowout or endpoint is inferred from positive geometric scales.


\section{Source objects and the finite calculation}
This calculation uses the preserved 165-page supplied manuscript
\emph{Finite Time Blowup for Navier--Stokes}, source equations
(7.23)--(7.42), (9.12)--(9.14). Its PDF SHA-256 is
\begin{center}\small
\texttt{8c8a94ad9ac824c8b605b9827cadf7beaca48bd10b380de3cfc872a2c37afa81}.
\end{center}
The source pulse fields, auxiliary rectangles, phases, and leading
covariance columns are the given objects; their global construction is
source input. Every finite map asserted below is proved here. The existing
coupled-stage proofs remain in the complete 102-page reader. The companion
pulse note proves the exact frame residual and a constructed finite-interval
intertwiner along its evolving coefficients. The source uses positive
$q(z,t)$, band scales $Q_\beta=2^{-\ell_\beta}$,
$\varepsilon_\beta=Q_\beta^h$, $A=1/2+h$; none is set to one here or
identified with the original coupled $A_0,\lambda_0,\mu,\nu=\sqrt{A_0}$
or its physical diffusivities.

All operations in this note are on a compact set inside $q>0$ and the
source active annulus. There are finitely many relevant bands and
locally finite source labels. Supports of newly prescribed inputs are
strictly inside the annulus, so no lower bound on a leading amplitude
at a vanishing shell edge is assumed. Endpoint and shell-edge estimates
are not inferred from compact-interior estimates.

\section{Physical curl of the actual retained amplitude}
Let $y=(\Theta,\Omega)^T$ be the actual smooth retained coupled solution,
$y'=q_v(C_cy+c_c)$, with the coefficients of the companion pulse note.
Use the smooth actual coefficients on the chosen slow/auxiliary patch;
all harmonic coefficients and cutoffs are independent of $\theta$ in
cylindrical components, and each carrier has nonzero integer angular
frequency. For its source tangent frame $B$, smooth coefficient map $L$,
and a fixed smooth real compact cutoff $\chi$ of this type, put
\[
 t_m=\chi BLy,\qquad
 \mathcal E_\chi=
 \chi\{[L'+q_vLC_c-HL+m^2d_NL]y+q_vLc_c\}+\chi'Ly,
 \quad H=B_\ell(\mathsf A_\Phi B-B').
\]
Here prime differentiates along the pulse coordinate; the other slow
and auxiliary derivatives remain in the spatial formulas below.
Smooth parameter dependence of the companion's constructed $L$ follows
directly from its integral equations on a fixed interval: for a nonzero
parameter multi-index $I$,
\[
 \partial^I X(v)=
 \int_{v_0}^v\left[D\,\partial^I X+
 \sum_{0<J\le I}{I\choose J}(\partial^JD)
                                 \partial^{I-J}X\right]\,\dd w .
\]
The initial parameter derivatives vanish for identity initial data.
Inductively the known forcing is smooth and bounded on the compact patch;
the same convergent ordered-integral iteration solves this equation.
This proves every finite mixed derivative of $X,X^{-1},L$.
Since $n^TB=0$, $n^Tt_m=0$. Differentiation proves
$t_m'-\mathsf A_\Phi t_m+m^2d_Nt_m=B\mathcal E_\chi$.
Thus even the pulse-cutoff term is explicit.

A physical version of the exact curl construction makes the intermediate
map to momentum residual unambiguous. For each chosen harmonic let
$\phi(x,t)=k_\beta m\Phi_\beta(x,t,Y(x,t))$ be its real evaluated phase,
let $\xi=\nabla_x\phi\ne0$, and let $a(x,t)=Q_\beta^{-A}t_m$ be the
physical transverse amplitude. The source coordinate change is
$r=\sqrt{Q_\beta}R$, $z=Q_\beta^D Z$, $1-t=Q_\beta T$,
with $D=1/2-h$ and normalized axial derivative
$D_z=\varepsilon_\beta\partial_Z$. It gives
$\xi=(k_\beta m/\sqrt{Q_\beta})n_\Phi$; all auxiliary chain derivatives
are included in $n_\Phi$. Hence $\xi\cdot a=0$. Define in Cartesian
components
\begin{equation}
 C=\frac{i\,\xi\times a}{|\xi|^2},\quad
 \mathcal A=C e^{i\phi},\quad
 \nabla\times\mathcal A=(a+r_a)e^{i\phi},\qquad
 r_a=\nabla\times C. \label{nsbridge:cb:curl}
\end{equation}
The vector triple-product identity gives
$i\xi\times C=a$, proving the last equality. It also shows that the
physical potential has exactly the source factor $Q_\beta^{1/2-A}$
relative to its chart coefficient
$i\,n_\Phi\times t_m/(k_\beta m|n_\Phi|^2)$.
Cartesian differentiation of $C$ includes derivatives of the cylindrical
basis; the same terms are the $R^{-1}$ terms of source (7.37)--(7.38).
Take real parts and sum the finite retained harmonics. This constructs
a smooth compact velocity $w_y$ with $\nabla\cdot w_y=0$.

Retain a finite smooth source state $(u_*,p_*)$, including its background,
actual leading waves and their curl corrections, and the real
physical pressure increments obtained from the companion pulse pressure
formula with $f_m=0$:
\[
 \pi_{\chi,m}=\frac{i}{k_\beta m|n|^2}
                  \bigl(n^T\mathsf Kt_m-(n')^Tt_m\bigr),\qquad
 \pi_y=\sum_{\beta,m}\operatorname{Re}
             (Q_\beta^{-2A}\pi_{\chi,m}e^{i\phi_{\beta,m}}).
\]
These are smooth and compact wherever $t_m$ is. This definition fixes the
pressure map without assuming the pulse residual vanishes.
Define the actual candidate residual
\begin{align}
 R_y={}&R_{\nu_{\rm NS}}(u_*+w_y,p_*+\pi_y), \nonumber\\
 R_{\nu_{\rm NS}}(u,p)
   :={}&\partial_tu+(u\cdot\nabla)u-\nu_{\rm NS}\Delta u+\nabla p .
 \label{nsbridge:cb:Ry}
\end{align}
This is a definition by explicit differentiated fields, not an assumed
source term. Its exact expansion is
\[
 R_y=R_{\nu_{\rm NS}}(u_*,p_*)
 +\partial_tw_y+(u_*\cdot\nabla)w_y+(w_y\cdot\nabla)u_*
 -\nu_{\rm NS}\Delta w_y+\nabla\pi_y
 +(w_y\cdot\nabla)w_y .
\]
Every cutoff, parent, phase, and curl derivative occurs in this formula.
Physical $\nu_{\rm NS}$ is retained and is distinct from every source
time-scaling convention. Define
$\Delta R_y=R_y-R_{\nu_{\rm NS}}(u_*,p_*)$.
We correct the new compact increment; this difference has compact support
since every added term does. The baseline residual remains explicitly.

The averaging needed next is the source extended-field operation, not a
spatial average of the evaluated field. Evaluate derivatives on the
extended coefficients first by the complete chain rule. In a fixed
chart set $\mathcal D_i=\partial_{x_i}
+\sum_\alpha(\partial_{x_i}Y_\alpha)\partial_{Y_\alpha}$ and similarly
$\mathcal D_t$. These are the pullbacks of commuting physical
derivatives, so $[\mathcal D_i,\mathcal D_j]=0$ on their lifted domain.
For smooth periodic coefficients their auxiliary derivative terms have
zero Haar average. In cylindrical components angular differentiation
also differentiates the basis. This defines an extended residual
$\widetilde R_y$ whose evaluation is exactly \eqref{nsbridge:cb:Ry}.
Average the tangential components of the increment at fixed slow variables:
\begin{equation}
 F_2=\av{\Delta\widetilde R_{y,\theta}},\qquad
 F_1=\av{\Delta\widetilde R_{y,z}},\qquad
 \av{g}=\int_{\mathbb T^2}\frac1{2\pi}
                    \int_0^{2\pi}g\,\dd\theta\,\dd Y . \label{nsbridge:cb:F}
\end{equation}
These $F_e$ are smooth and compact
in the chosen physical annulus. Equations \eqref{nsbridge:cb:curl}--\eqref{nsbridge:cb:F}
are the explicit intermediate map from the retained coupled pair to a
signed momentum residual. They do not identify it with the original
two-dimensional Boussinesq force.

\section{Radial moments, support, and the original chart scales}
For each fixed $(z,t)$, choose $0<r_-<r_+$ containing the compact
supports of both $F_e$ and the chosen bumps in the active annulus.
The source bump is
\[
 b_e(r,z,t)=q^{-(e+1)/2}\widehat b_e(r/\sqrt q),\qquad
 \int_0^\infty x^e\widehat b_e(x)\,\dd x=1,\qquad e=1,2,
\]
with interior support. The change $r=\sqrt q\,x$ proves
$\int r^eb_e\,\dd r=1$ exactly. Define
\begin{equation}
 M_e=\int_0^\infty r^eF_e(r)\,\dd r,\qquad
 \sigma_e(r)=-r^{-e}\int_0^r s^e(F_e(s)-b_e(s)M_e)\,\dd s .
 \label{nsbridge:cb:primitive}
\end{equation}
Then
\begin{equation}
 (\partial_r+e/r)\sigma_e=-F_e+b_eM_e,\qquad
 \int_0^\infty r^e(F_e-b_eM_e)\,\dd r=0. \label{nsbridge:cb:radial}
\end{equation}
Differentiate $r^e\sigma_e$ to prove the first identity; integrate and
use the bump normalization for the second. The integral is zero below
$r_-$ and above $r_+$, so $\sigma_e$ has compact smooth support in the
same enclosing interval. A homogeneous solution is $cr^{-e}$, and a
compact homogeneous solution has $c=0$; thus the primitive is unique
among compact solutions of its displayed equation. Moreover a compact
solution of $(\partial_r+e/r)\sigma=-F_e$ exists if and only if $M_e=0$:
necessity follows by integrating the weighted derivative, and
sufficiency follows from \eqref{nsbridge:cb:primitive}. The two moments are
therefore the exact surviving finite-dimensional terms.

For $C_e=\int_{r_-}^{r_+}s^e\,\dd s$,
$|M_e|\le C_e\|F_e\|_\infty$ and
\[
 \|\sigma_e\|_\infty\le r_-^{-e}C_e
        (\|F_e\|_\infty+\|b_e\|_\infty|M_e|).
\]
If $G_e=F_e-b_eM_e$, successive radial derivatives obey the exact recurrence
\[
 \partial_r^{j+1}\sigma_e=-\partial_r^jG_e
 -e\sum_{i=0}^j {j\choose i}(-1)^ii!r^{-i-1}
                              \partial_r^{j-i}\sigma_e .
\]
This proves every finite radial derivative bound inductively. Parameter
derivatives follow by differentiating the compact integral; all derivatives
of $q,b_e$ and the moving physical data are retained. On the present compact
set $q$ is bounded away from zero. No terminal uniform estimate is implied.

The source (9.12)--(9.14) uses physical
$F_e=Q^{-2A-1/2}E_e^*$, where
$E_2^*=\av{E_\theta}_Y$ and $E_1^*=\av{E_z}_Y$.
For our fields define $E_e^*=Q^{2A+1/2}F_e$ in that same convention.
With $r=\sqrt Q R$ and $Q$ fixed in a chart, the exact conversions are
\begin{gather}
 M_e=Q^{e/2-2A}M_e^*,\quad M_e^*=\int R^eE_e^*\,\dd R,\quad
 b_e^*(R)=Q^{(e+1)/2}b_e(\sqrt Q R),\nonumber\\
 \Sigma_e(R)=Q^{2A}\sigma_e(\sqrt Q R),\qquad
 (\partial_R+e/R)\Sigma_e=-E_e^*+b_e^*M_e^*. \label{nsbridge:cb:chart}
\end{gather}
Substitute $r=\sqrt Q R$ including its measure in the moment, and
differentiate $\Sigma_e$ to prove each factor. The input $F_e$ is already
auxiliary independent. Taking only an averaged moment of an unaveraged
$F(R,Y)$ would not suffice: $F=b_e(R)\cos(2\pi Y_1)$ has zero averaged
moment but a nonzero exterior primitive at each generic $Y$.

\section{The source covariance derivative and its signed right inverse}
At each slow point let $w_{\rm tan}=(w_\theta,w_z)$ and define
\[
 \Cov(w)=\av{w_rw_{\rm tan}},\qquad
 \Cross(w,v)=\av{w_rv_{\rm tan}+v_rw_{\rm tan}}.
\]
Direct expansion proves
$\Cov(w+v)=\Cov(w)+\Cross(w,v)+\Cov(v)$ and
$D\Cov(w)[v]=\Cross(w,v)$.

Use the exact source real pulses
$b_\sigma=\chi_g(\xi_g)\psi(v)t_\sigma^h\cos(k\Phi_\sigma)$,
$\sigma\in\{+,-\}$. They exclude the slow cutoff $\eta_\beta$ and
scalar amplitudes; $\chi_g$ is an auxiliary transverse cutoff.
The two signs have disjoint auxiliary supports. Put
$H_{\rm cov}=[\Cov(b_+)\mid\Cov(b_-)]$.
For slow scalar coefficients independent of both averages,
\[
 \Cov(a_+b_++a_-b_-)=H_{\rm cov}(a_+^2,a_-^2)^T.
\]
All cross products vanish by disjointness. The factor $1/2$ from
$\av{\cos^2(k\Phi_\sigma)}_\theta$ is already inside the columns.

On the source interior cone the fixed primary amplitudes satisfy
$a_\sigma=\sqrt{(H_{\rm cov}^{-1}T_{0,*})_\sigma}>0$,
$T_{0,*}=(Q/q)^{A+1/2}T_0$, and
$W_0=\sqrt\varepsilon\sum_\sigma a_\sigma b_\sigma$.
For any signed slow real vector $\Sigma$, define
\begin{equation}
 d_\Sigma=H_{\rm cov}^{-1}\Sigma/\varepsilon,\qquad
 \delta a_\sigma=(d_\Sigma)_\sigma/(2a_\sigma),\qquad
 L_\Sigma=\sqrt\varepsilon\sum_\sigma\delta a_\sigma b_\sigma .
 \label{nsbridge:cb:L}
\end{equation}
It follows exactly that
\begin{gather}
 \Cross(W_0,L_\Sigma)=\varepsilon H_{\rm cov}(2a\odot\delta a)=\Sigma,
 \label{nsbridge:cb:inverse}\\
 \Cov(W_0+L_\Sigma)
   =\varepsilon T_{0,*}+\Sigma+
       \varepsilon H_{\rm cov}(\delta a_+^2,\delta a_-^2)^T.
 \label{nsbridge:cb:quad}
\end{gather}
This proves a linear right inverse of the derivative. It is not an
exact inverse of the quadratic covariance map. There is no sign or
size restriction on $\Sigma$ in these identities; no square root of
$H_{\rm cov}^{-1}(T_{0,*}+\Sigma/\varepsilon)$ is taken.

With Euclidean and induced operator norms and
$a_{\min}=\min_\sigma|a_\sigma|>0$, they give the explicit bound
\begin{equation}
 \|\Cov(L_\Sigma)\|\le
 \frac{\|H_{\rm cov}\|\|H_{\rm cov}^{-1}\|^2}
              {4\varepsilon a_{\min}^2}\,\|\Sigma\|^2. \label{nsbridge:cb:bound}
\end{equation}
Indeed $\|\delta a\|\le\|H_{\rm cov}^{-1}\|\|\Sigma\|/
(2\varepsilon a_{\min})$ and
$\|(\delta a_\sigma^2)_\sigma\|_2\le\|\delta a\|_2^2$.
All higher derivative costs are explicit from Leibniz and
\[
 \partial^I H_{\rm cov}^{-1}
 =-H_{\rm cov}^{-1}
 \sum_{0<J\le I}{I\choose J}(\partial^JH_{\rm cov})
                           \partial^{I-J}H_{\rm cov}^{-1},
\]
and, for any nonzero scalar $a$,
$\partial^I(a^{-1})=-a^{-1}\sum_{0<J\le I}{I\choose J}
(\partial^Ja)\partial^{I-J}(a^{-1})$.
Together with \eqref{nsbridge:cb:L} these recurrences give finite bounds for
each order without assuming constant amplitudes or constant columns.

\section{Assembly, the actual covariance change, and residual cancellation}
For the physical target $\sigma=(\sigma_2,\sigma_1)$ from
\eqref{nsbridge:cb:primitive}, use the source squared slow partition
$\sum_\beta\eta_\beta^2=1$ on its support. Set
\begin{align}
 W_{0,\rm phys}&=\sum_\beta Q_\beta^{-A}\eta_\beta W_{0,\beta},\nonumber\\
 v_{\rm pr}&=\sum_\beta Q_\beta^{-A}\eta_\beta
                        L_\beta(Q_\beta^{2A}\sigma).
 \label{nsbridge:cb:assembly}
\end{align}
Overlapping slow labels have disjoint auxiliary supports. Taking cross
covariance and using \eqref{nsbridge:cb:inverse} gives
\[
 \Cross(W_{0,\rm phys},v_{\rm pr})
 =\sum_\beta Q_\beta^{-2A}\eta_\beta^2Q_\beta^{2A}\sigma=\sigma.
\]
Every band scale remains distinct. Likewise the leading physical factor is
$Q^{-2A}\varepsilon(Q/q)^{A+1/2}
=Q^{-A+h+1/2}q^{-A-1/2}=q^{-A-1/2}$ because $A=1/2+h$.

Apply the exact curl \eqref{nsbridge:cb:curl} to each coefficient of
$v_{\rm pr}$, including the slow cutoffs before curling. Write the
resulting divergence-free increment as $v_s=v_{\rm pr}+r_s$.
Let $t_{s,m}$ denote the chart transverse harmonic coefficients of
$\eta_\beta L_\beta(Q_\beta^{2A}\sigma)$, before multiplication by
$Q_\beta^{-A}$. For each such coefficient choose
$\pi_{s,m}=i(n^T\mathsf Kt_{s,m}-(n')^Tt_{s,m})/
(k_\beta m|n|^2)$ and sum
$\operatorname{Re}(Q_\beta^{-2A}\pi_{s,m}e^{i\phi_{\beta,m}})$
to define the smooth compact $\pi_s$. Its full gradients remain below.
Let the actual old wave $w$ consist of the wave part of $u_*$
plus the candidate $w_y$, including all existing curl corrections, and write
$w=W_{0,\rm phys}+e_w$. Define tensor covariance
$\mathcal T(w)=\av{w\otimes w}$ and
$\mathcal B_T(u,v)=\av{u\otimes v+v\otimes u}$.
Expansion, with $e_w=w-W_{0,\rm phys}$, proves
\begin{equation}
 \Delta\mathcal T=
 \mathcal B_T(W_{0,\rm phys},v_{\rm pr})
 +\mathcal B_T(e_w,v_{\rm pr})
 +\mathcal B_T(w,r_s)+\mathcal T(v_s).
 \label{nsbridge:cb:fullcov}
\end{equation}
Its first radial--tangential pair is $\sigma$.
All other entries and the remaining three tensors are retained.
For example their radial--tangential pair has bound
\[
 2\|e_w\|_{L^2_{\theta,Y}}\|v_{\rm pr}\|_{L^2_{\theta,Y}}
 +2\|w\|_{L^2_{\theta,Y}}\|r_s\|_{L^2_{\theta,Y}}
 +\|v_s\|_{L^2_{\theta,Y}}^2 ,
\]
by Cauchy--Schwarz. The compact-set estimates above and the explicit
curl differentiate to bound all these terms at each finite order.

For clarity the complete residual after adding $v_s$ and its selected
pressure $\pi_s$ is, without any averaging,
\begin{align}
 R_{\nu_{\rm NS}}(u+v_s,p+\pi_s)
 =R_{\nu_{\rm NS}}(u,p)
 &+\partial_tv_s+(u\cdot\nabla)v_s+(v_s\cdot\nabla)u \nonumber\\
 &-\nu_{\rm NS}\Delta v_s+\nabla\pi_s
 +(v_s\cdot\nabla)v_s . \label{nsbridge:cb:residual}
\end{align}
This follows by expanding the product and retains viscosity and every
pressure term. For a symmetric tensor $S$ the averaged tangential
divergence in physical cylindrical coordinates is
\[
 (\av{\nabla\cdot S})_\theta
   =(\partial_r+2/r)\av{S_{r\theta}}+\partial_z\av{S_{z\theta}},
 \qquad
 (\av{\nabla\cdot S})_z
   =(\partial_r+1/r)\av{S_{rz}}+\partial_z\av{S_{zz}} .
\]
To verify, write the Cartesian divergence in the orthonormal cylindrical
basis, using $\partial_\theta e_r=e_\theta$,
$\partial_\theta e_\theta=-e_r$; periodic angular derivatives average
to zero and the two basis terms give $2/r$ and $1/r$ respectively.
The auxiliary chain terms average to zero by periodicity, but the
physical slow derivatives remain.

Let $U=u_*+w_y$, $P_y=p_*+\pi_y$, and set
\[
 \mathcal L=
 \partial_tv_s+(U\cdot\nabla)v_s+(v_s\cdot\nabla)U
 -\nu_{\rm NS}\Delta v_s+\nabla\pi_s .
\]
For a fully explicit averaged remainder, write $U=u_{\rm mean}+w$,
where $u_{\rm mean}$ is its angular mean, and put
$\mathcal V=\mathcal L-(w\cdot\nabla)v_s-(v_s\cdot\nabla)w$.
Then the average of \eqref{nsbridge:cb:residual} is
$\av{R_y}+\av{\mathcal V}+\nabla\cdot\Delta\mathcal T$.
Let $S_0=\mathcal B_T(W_{0,\rm phys},v_{\rm pr})$ and
$S_{\rm rem}=\mathcal B_T(e_w,v_{\rm pr})
 +\mathcal B_T(w,r_s)+\mathcal T(v_s)$.
Consequently the radial divergence of the first pair in
\eqref{nsbridge:cb:fullcov} changes the prescribed averaged residual $F_e$ into
$b_eM_e$ by \eqref{nsbridge:cb:radial}. The axial tensor divergences, the other
three tensors, the exact linear terms and the pressure terms of
\eqref{nsbridge:cb:residual} remain in the new residual. Equivalently the new
averaged tangential residual is exactly
\begin{align*}
 \av{R_{\rm new}}_{\rm tan}
 ={}&\av{R_{\nu_{\rm NS}}(u_*,p_*)}_{\rm tan}
 +(b_2M_2,b_1M_1)+\av{\mathcal V}_{\rm tan}\\
 &+\partial_z((S_0)_{z\theta},(S_0)_{zz})
   +(\nabla\cdot S_{\rm rem})_{\rm tan}.
\end{align*}
Every term here is defined by the displayed differentiated fields and
tensors. None is assumed to vanish.
The map from $y$ through \eqref{nsbridge:cb:F}, \eqref{nsbridge:cb:primitive},
\eqref{nsbridge:cb:L}, and \eqref{nsbridge:cb:assembly} is now an actual composite
construction. It retains the signed residual produced by the coupled
candidate and the source's original physical and chart units.

\section{Signs and the exact full-flow reflection}
The maps $F_e\mapsto M_e\mapsto\sigma_e\mapsto L_\Sigma$ are linear
and reverse sign with their inputs; $\Cov(L_\Sigma)$ is even.
The two primary amplitudes may also be replaced independently by
$s_\sigma a_\sigma$, $s_\sigma\in\{-1,1\}$, provided the same fixed
choice is used in each denominator: the cross identity is unchanged.
This does not negate the nonlinear Navier--Stokes equation.

Here is the exact full-flow map. In Cartesian coordinates let
$O=\operatorname{diag}(1,-1,1)$ and define
\[
 u^\sharp(x,t)=O u(Ox,t),\quad
 p^\sharp(x,t)=p(Ox,t),\quad f^\sharp(x,t)=O f(Ox,t).
\]
The chain rule gives
$\nabla u^\sharp=O(\nabla u)(Ox,t)O$,
$\Delta u^\sharp=O(\Delta u)(Ox,t)$, and
$(u^\sharp\cdot\nabla)u^\sharp
 =O((u\cdot\nabla)u)(Ox,t)$.
Thus $\nabla\cdot u^\sharp=0$ and
$R_{\nu_{\rm NS}}(u^\sharp,p^\sharp)=O R_{\nu_{\rm NS}}(u,p)(Ox,t)$
with the same positive viscosity.
In cylindrical components this is
\[
 (u_r^\sharp,u_\theta^\sharp,u_z^\sharp)(r,\theta,z,t)
 =(u_r,-u_\theta,u_z)(r,-\theta,z,t),
\]
and the force transforms identically. Pressure reflects in its argument.
For an axisymmetric field this reverses exactly its swirl; meridional
components remain fixed. The entire momentum tensor is mapped to
$O\mathcal T O^T$, so its radial--tangential pair transforms by
$(T_{r\theta},T_{rz})\mapsto(-T_{r\theta},T_{rz})$ after angular
averaging. Vorticity transforms as
$\omega^\sharp(x,t)=\det(O)\,O\omega(Ox,t)$, obtained by applying
the same derivative identities to the Levi-Civita formula for curl.
The map is an involution, preserves the kinetic-energy integral
(by $|\det O|=1$), compact support, smoothness, and blowup magnitude
along the reflected path. It preserves zero initial data when present.
It proves transport of any established axisymmetric signed swirl
divergence to the opposite sign; it does not itself prove an endpoint.

The source five-moment correction, auxiliary mean inverse, and infinite
summation are subsequent source operations. Their full estimates are
not re-proved here. The earlier draft formulas for them are withdrawn:
in particular the fifth moment row is
$\int(2RG\gamma_d-RV\Delta v)\,\dd R=-J_z$, without an extra outer
$R$, and its reserved profile in a $Q$ chart includes $(Q/q)^A$.
The finite map proved above retains the two moments and all other
residual terms, so it establishes neither a third shooting return nor
an infinite modified viscous sequence.
